\section{Introducción: de AlphaGo a AlphaProof}
Thomas Hubert ubica esta clase en la unidad de «razonamiento-ejecución» de Berkeley LLM Agents SP25, y usa AlphaZero/AlphaTensor para enlazar la época dorada del aprendizaje por refuerzo con AlphaProof. Hubert señala explícitamente que la \textbf{rigurosidad}, la \textbf{feedback demostrable} y la \textbf{creatividad ilimitada} que exige el razonamiento matemático son altamente isomorfas a la \textbf{búsqueda}, la \textbf{generalización} y el \textbf{retorno confiable} que exige la exploración de fronteras en RL. Esto explica por qué DeepMind sitúa a las matemáticas como «el siguiente campo de batalla de RL».

\begin{knowledgebox}{La resonancia entre las matemáticas y RL}
Las matemáticas son una red de alta dimensión donde los nodos son teoremas y las aristas son cadenas lógicas; RL, por su parte, explora en el espacio de acciones. Hubert enfatiza: «en cuanto puedas abstraer un teorema como una política de acciones, el optimizador puede encargarse de aprender la teoría.»
\end{knowledgebox}

La clase se divide en tres partes: la primera repasa la evolución de la formalización matemática y organiza Lean/Mathlib; la segunda resume el ciclo virtuoso de AlphaZero y cómo AlphaProof se conecta con los sistemas de demostración; la tercera se centra en la práctica de ingeniería, la reproducción de la IMO y el ciclo de gobernanza. Toda la clase sigue la lógica pedagógica de «teoría → plataforma → aplicación», y en cada parte se reserva tiempo para mostrar lecciones aprendidas y preguntas y respuestas.

\begin{figure}[H]
\centering
\includegraphics[width=0.8\textwidth,height=0.35\textheight,keepaspectratio]{\videocoverpath}
\caption{Portada de la Lecture 05 y esquema del flujo de AlphaProof (la portada aporta la imagen principal y una imagen arquitectónica, y resalta la fusión de RL con la formalización).}
\end{figure}

Hubert menciona en particular que esta clase es la quinta de SP25, e inmediatamente después de «agentes generativos y razonamiento». Considera a AlphaProof como «un agente en la versión matemática», que necesita combinar el retorno determinista con la gobernanza interactiva para que la búsqueda de demostraciones de alta confiabilidad sea posible.

\subsection{Resumen de la sección}
La parte introductoria se encarga de construir el contexto: la similitud estructural entre las matemáticas y RL, la misión de AlphaProof, y la ubicación en el ritmo de toda la serie de clases de SP25 entre la «cadena de razonamiento y la práctica de ingeniería», sentando así la ruta causal para el contenido posterior.
\newpage

